\documentclass[11pt,oneside]{report}

\usepackage[margin=1in]{geometry}
\usepackage[T1]{fontenc}
\usepackage[utf8]{inputenc}
\usepackage{amsmath}
\usepackage{amssymb}
\usepackage{booktabs}
\usepackage{longtable}
\usepackage{array}
\usepackage{enumitem}
\usepackage{xcolor}
\usepackage{listings}
\usepackage{hyperref}

\hypersetup{
  colorlinks=true,
  linkcolor=blue,
  urlcolor=blue,
  citecolor=blue,
  pdftitle={Auto Aim Onboarding Course},
  pdfauthor={auto aim maintainers}
}

\lstset{
  basicstyle=\ttfamily\small,
  breaklines=true,
  frame=single,
  columns=fullflexible,
  keepspaces=true,
  showstringspaces=false,
  tabsize=2
}

\setlist[itemize]{topsep=4pt,itemsep=2pt}
\setlist[enumerate]{topsep=4pt,itemsep=2pt}

\DeclareRobustCommand{\repo}{\texttt{auto\_aim}}
\DeclareRobustCommand{\code}[1]{\texttt{\detokenize{#1}}}
\DeclareRobustCommand{\topic}[1]{\texttt{\detokenize{#1}}}
\DeclareRobustCommand{\param}[1]{\texttt{\detokenize{#1}}}
\DeclareRobustCommand{\file}[1]{\texttt{\detokenize{#1}}}
\newcommand{\vect}[1]{\mathbf{#1}}
\pdfstringdefDisableCommands{%
  \def\repo{auto aim}%
  \def\code#1{#1}%
  \def\topic#1{#1}%
  \def\param#1{#1}%
  \def\file#1{#1}%
}

\title{
  Auto Aim Onboarding Course\\
  \large Mathematics, Coding, and Source Walkthrough
}
\author{RoboMaster auto\_aim package}
\date{May 10, 2026}

\begin{document}

\maketitle
\tableofcontents

\chapter*{How To Use This Course}
\addcontentsline{toc}{chapter}{How To Use This Course}

This document is written for new members who may have no previous coding
experience. It is not only a reference manual. It is meant to be read as a
course:

\begin{enumerate}
  \item Read the mathematics course first. The code will make more sense once
  the coordinate frames, prediction, and filtering ideas are familiar.
  \item Read the coding course next. It introduces C++, ROS 2, and the build
  system using this repository as the example.
  \item Use the code walkthrough while opening the source files in an editor.
  Read one file at a time and compare the prose with the implementation.
  \item Debug with the structured topic \topic{/auto_aim/debug}. Most failures
  can be located by following the same order as the pipeline.
\end{enumerate}

The most important rule for this repository is simple: do not tune a downstream
parameter to hide an upstream error. For example, do not change
\param{pitch_offset_deg} to compensate for an unmeasured bullet speed or an
incorrect barrel offset. Fix the earliest wrong stage first.

To compile this document from the workspace root, run:

\begin{lstlisting}
cd src/docs/latex
pdflatex main.tex
pdflatex main.tex
\end{lstlisting}

Running \code{pdflatex} twice updates the table of contents.

\chapter*{Repository Map}
\addcontentsline{toc}{chapter}{Repository Map}

\begin{longtable}{p{0.32\textwidth}p{0.60\textwidth}}
\toprule
Path & Purpose \\
\midrule
\file{src/src/auto_aim_node.cpp} & The ROS 2 node. Owns subscriptions,
publishers, parameters, callbacks, command smoothing, and debug publication. \\
\file{src/src/pnp_solver.cpp} & Converts a 2D detector bounding box into a 3D
pose using OpenCV PnP. \\
\file{src/src/frame_transformer.cpp} & Converts camera-frame PnP output into
the odom/microcontroller reference frame. \\
\file{src/src/tracker.cpp} & Extended Kalman filter, target association,
state machine, face prediction, and ballistic aim planning. \\
\file{src/src/fire_gate.cpp} & Central fire/hold decision logic and blocker
reason generation. \\
\file{src/src/detection_adapter.cpp} & Converts raw detector messages into
small validated detection candidates. \\
\file{src/src/debug_publisher.cpp} & Copies per-frame debug data into the
\topic{/auto_aim/debug} ROS message and logs blocker histograms. \\
\file{src/src/config_validator.cpp} & Checks startup parameters for impossible
or suspicious values. \\
\file{src/include/auto_aim/*.hpp} & Public declarations for the C++ modules. \\
\file{src/msg/AutoAimDebug.msg} & Structured debug message for one frame. \\
\file{src/msg/GimbalCmd.msg} & Placeholder message, not used at runtime. \\
\file{src/config/*.yaml} & Robot-specific parameter files. \\
\file{src/launch/*.py} & ROS 2 launch files. \\
\file{src/docs/*.md} & Existing operational guides. \\
\bottomrule
\end{longtable}

\chapter{Mathematics Course}

\section{What The Math Is For}

The auto-aim system answers one question many times per second:

\begin{quote}
Given a camera image, the current gimbal yaw and pitch, and a moving target,
what absolute yaw and pitch should the gimbal move to, and is it safe to fire?
\end{quote}

The repository solves that question as a chain of smaller math problems:

\begin{enumerate}
  \item A detector reports a 2D rectangle around an armor plate.
  \item PnP estimates where that plate is in 3D camera coordinates.
  \item A frame transform moves the 3D point into the robot's odom frame.
  \item An Extended Kalman Filter estimates the enemy robot center, velocity,
  yaw, yaw rate, and armor radius.
  \item The aim planner predicts which of four armor faces should be hit at
  bullet impact time.
  \item The ballistic solver computes the pitch needed to compensate gravity.
  \item The fire gate checks whether the command is accurate enough to shoot.
\end{enumerate}

Each step has its own units. Mixing units is one of the fastest ways to break
this project. In this package, distances are meters unless a comment says
millimeters, angles are radians unless a parameter name ends with
\code{_deg}, and image coordinates are pixels.

\section{Math Notation For Beginners}

A scalar is one number, such as $3.0$ meters. A vector is a list of numbers,
such as a 3D position:

\[
\vect{p} =
\begin{bmatrix}
x \\ y \\ z
\end{bmatrix}.
\]

A matrix is a rectangular table of numbers. Multiplying a matrix by a vector
can rotate, scale, or mix the vector components. A rotation matrix $R$ and a
translation vector $\vect{t}$ transform a point from one frame to another:

\[
\vect{p}_{world} = R \vect{p}_{camera} + \vect{t}.
\]

The code often uses the same idea with quaternions instead of explicit
rotation matrices. A quaternion is another way to store a 3D rotation. You do
not need to manually multiply quaternions at first; you need to know that
\code{tf2::quatRotate(q, p)} rotates point $\vect{p}$ by rotation $q$.

\section{Coordinate Frames}

A coordinate frame defines what the axes mean. The same physical armor can
have different coordinates in different frames.

\begin{longtable}{p{0.24\textwidth}p{0.65\textwidth}}
\toprule
Frame & Meaning in this repository \\
\midrule
Camera optical & ROS camera convention: $+Z$ forward from the camera,
$+X$ right in the image, $+Y$ down in the image. PnP outputs points here. \\
Gimbal body & Mechanical convention: $+X$ forward, $+Y$ left, $+Z$ up. The
barrel offset is measured in this frame. \\
Odom / micro reference & Fixed startup reference used by the microcontroller.
The node publishes absolute yaw and pitch commands in this frame after sign
conversion. \\
\bottomrule
\end{longtable}

In \file{src/src/frame_transformer.cpp}, the current IMU yaw and pitch build a
rotation:

\[
q_{camera\rightarrow odom} = q_{gimbal} q_{body\rightarrow camera}.
\]

Then a PnP point is transformed as:

\[
\vect{p}_{odom} =
q_{camera\rightarrow odom} \vect{p}_{camera} +
\begin{bmatrix}0\\0\\h_{gimbal}\end{bmatrix}.
\]

The code form is:

\begin{lstlisting}[language=C++]
tf2::Vector3 p_odom = tf2::quatRotate(rotation_, p_cam) + translation_;
\end{lstlisting}

\subsection{Yaw And Pitch}

Yaw turns left/right around the vertical axis. Pitch turns up/down. The command
published on \topic{/tracker/cmd_gimbal} is not a small change. It is an
absolute destination:

\[
\text{command} =
\begin{cases}
\text{absolute yaw target in radians}\\
\text{absolute pitch target in radians}
\end{cases}
\]

The parameters \param{gimbal.yaw_sign} and \param{gimbal.pitch_sign} convert
between the code's internal convention and the microcontroller's convention.
They must be exactly $+1$ or $-1$.

\section{Camera Projection}

A camera maps a 3D point to a 2D pixel. Ignoring distortion for the moment:

\[
u = f_x \frac{X}{Z} + c_x,
\qquad
v = f_y \frac{Y}{Z} + c_y.
\]

Here $(X,Y,Z)$ is a point in camera optical coordinates, $(u,v)$ is the pixel,
$f_x$ and $f_y$ are focal lengths in pixels, and $(c_x,c_y)$ is the optical
center in pixels. These values come from \topic{/camera_info}.

The code also uses the inverse idea for the HUD overlay. If the aim ray has a
small yaw and pitch in the camera frame, the pixel is:

\[
u = c_x + f_x \tan(\theta_{yaw}),
\qquad
v = c_y - f_y \tan(\theta_{pitch}).
\]

This appears in \file{src/src/auto_aim_node.cpp} when publishing
\topic{/tracker/aim_pixels}.

\section{PnP: From 2D Rectangle To 3D Pose}

PnP means ``Perspective-n-Point''. It estimates the camera-frame pose of a
known 3D object from corresponding 3D object points and 2D image points.

This repository does not get exact armor corner points from the detector. It
gets a bounding box: center $(c_x,c_y)$, width $w$, and height $h$. The PnP
solver builds a synthetic rectangle inside that box:

\[
x_{left} = c_x - \rho \frac{w}{2},
\qquad
x_{right} = c_x + \rho \frac{w}{2},
\]

\[
y_{top} = c_y - \rho \frac{h}{2},
\qquad
y_{bottom} = c_y + \rho \frac{h}{2},
\]

where $\rho$ is \param{light_ratio}. The four image points are:

\[
(x_{left}, y_{bottom}),\quad
(x_{left}, y_{top}),\quad
(x_{right}, y_{top}),\quad
(x_{right}, y_{bottom}).
\]

The 3D object model is a flat armor plate. The solver tries both dimensions:

\begin{itemize}
  \item small armor: $140$ mm by $125$ mm,
  \item large armor: $235$ mm by $127$ mm.
\end{itemize}

It keeps the model with lower reprojection error. Reprojection error means:

\begin{enumerate}
  \item Estimate the pose.
  \item Project the known 3D corners back into the image.
  \item Measure the average pixel distance between projected corners and input
  corners.
\end{enumerate}

The code rejects a PnP result when the mean reprojection error is too large.
The debug topic also reports a normalized error:

\[
e_{norm} = \frac{e_{pixels}}{\sqrt{w^2+h^2}}.
\]

This makes errors easier to compare across targets of different pixel size.

\section{Armor Yaw And Obliquity}

The PnP solver estimates both position and orientation. Orientation can flip
for flat rectangles because two poses can project similarly into the image.
The frame transformer protects the tracker from a bad plate normal.

Let the transformed armor center be $(x,y,z)$ in odom. The bearing from the
robot to the armor is:

\[
\theta_{bearing} = \operatorname{atan2}(y,x).
\]

If the PnP yaw points the wrong way, the code adds $\pi$. If the yaw is still
too far from the bearing by more than \param{max_oblique_deg}, the code clamps
the yaw to the bearing. This means the position can still be used even when the
plate orientation is unreliable.

\section{Tracker State}

The tracker models the enemy robot as a rotating body with armor plates on a
circle around the robot center. The EKF state vector has nine values:

\[
\vect{x} =
\begin{bmatrix}
x_c & v_{xc} & y_c & v_{yc} & z_a & v_{za} & \psi & \dot{\psi} & r
\end{bmatrix}^T.
\]

\begin{longtable}{p{0.15\textwidth}p{0.65\textwidth}p{0.12\textwidth}}
\toprule
Symbol & Meaning & Unit \\
\midrule
$x_c,y_c$ & enemy robot center position in odom & m \\
$v_{xc},v_{yc}$ & enemy robot center horizontal velocity & m/s \\
$z_a$ & armor height & m \\
$v_{za}$ & armor vertical velocity & m/s \\
$\psi$ & visible armor face yaw & rad \\
$\dot{\psi}$ & armor yaw rate & rad/s \\
$r$ & radius from enemy center to armor plate & m \\
\bottomrule
\end{longtable}

The measurement vector is:

\[
\vect{z} =
\begin{bmatrix}
x_a & y_a & z_a & \psi
\end{bmatrix}^T,
\]

where $(x_a,y_a,z_a)$ is the detected armor center in odom.

\section{Measurement Model}

The tracker assumes the armor is offset from the enemy center by radius $r$:

\[
x_a = x_c - r\cos\psi,
\qquad
y_a = y_c - r\sin\psi,
\qquad
z_a = z_a,
\qquad
\psi = \psi.
\]

In EKF notation this is:

\[
\vect{z}_{pred} = h(\vect{x}).
\]

The innovation is the difference between what was measured and what the
current state predicted:

\[
\vect{y} = \vect{z} - h(\vect{x}).
\]

A small innovation means the new detection agrees with the tracker. A large
innovation means the detection may be noisy, a different target, or a different
armor face after the enemy rotated.

\section{Why The Tracker Is An EKF}

A normal Kalman filter works when the equations are linear. This tracker has
$\sin\psi$ and $\cos\psi$, so the measurement model is nonlinear. An Extended
Kalman Filter solves this by linearizing the nonlinear equation near the
current estimate.

The linearization is the Jacobian matrix $H$:

\[
H =
\frac{\partial h}{\partial \vect{x}}.
\]

For the first two measurement rows:

\[
\frac{\partial x_a}{\partial x_c}=1,\quad
\frac{\partial x_a}{\partial \psi}=r\sin\psi,\quad
\frac{\partial x_a}{\partial r}=-\cos\psi,
\]

\[
\frac{\partial y_a}{\partial y_c}=1,\quad
\frac{\partial y_a}{\partial \psi}=-r\cos\psi,\quad
\frac{\partial y_a}{\partial r}=-\sin\psi.
\]

These are exactly the non-zero entries built in
\file{src/src/tracker.cpp} inside \code{Tracker::ekfUpdate}.

\section{Prediction Step}

Between frames, the target is predicted with a damped constant-velocity model.
For a position and velocity pair:

\[
p_{new} = p + b v \Delta t,
\qquad
v_{new} = b v.
\]

The damping coefficient $b$ is computed from \param{alpha_pos},
\param{alpha_yaw}, \param{alpha_coast}, the frame time $\Delta t$, and
\param{ref_freq}. This keeps damping approximately consistent even if the
detector rate changes.

The covariance matrix $P$ stores uncertainty. During prediction it becomes:

\[
P_{new} = FPF^T + Q.
\]

$F$ is the Jacobian of the motion model. $Q$ is process noise. Larger
\param{q_pos} or \param{q_yaw} makes the tracker respond faster but jitter
more. Smaller values make it smoother but slower.

\section{Update Step}

When a detection is associated with the current track, the Kalman update is:

\[
S = HPH^T + R,
\]

\[
K = PH^T S^{-1},
\]

\[
\vect{x}_{new} = \vect{x} + K\vect{y}.
\]

$R$ is measurement noise. In this repository, $R$ grows with target distance.
The code also increases yaw and position noise when the armor is viewed
obliquely, because a nearly edge-on plate gives weaker orientation information.

The covariance update uses Joseph form:

\[
P_{new} = (I-KH)P(I-KH)^T + KRK^T.
\]

Joseph form costs a little more computation but is more numerically stable.

\section{Mahalanobis Distance}

The tracker must decide whether a detection belongs to the current target.
It uses Mahalanobis distance:

\[
d^2 = \vect{y}^T S^{-1} \vect{y}.
\]

This is better than plain Euclidean distance because it accounts for
uncertainty. A detection far from a very uncertain prediction may still be
reasonable. A detection close to a very certain prediction can be suspicious if
the covariance says it should not move there.

The parameter \param{maha_threshold} is the gate. If the distance is above
this threshold, the detection is not associated with the current track.

\section{Tracker State Machine}

The tracker has four states:

\begin{longtable}{p{0.22\textwidth}p{0.68\textwidth}}
\toprule
State & Meaning \\
\midrule
\code{LOST} & No target is trusted. The next valid detection initializes a new
track. \\
\code{DETECTING} & A possible target has been seen but needs consecutive matches
before commands are trusted. \\
\code{TRACKING} & The target is locked. The node may publish nonzero yaw/pitch
and the fire gate may allow shooting. \\
\code{TEMP_LOST} & The tracker had a target but missed recent detections. It
coasts briefly before becoming \code{LOST}. \\
\bottomrule
\end{longtable}

The purpose of this state machine is to avoid commanding the gimbal from a
single accidental detection.

\section{Armor Face Switching}

When an enemy spins, the visible armor face changes. The tracker detects a
large yaw jump. A jump near $90^\circ$ means the visible face changed to the
alternate pair. The code swaps \code{radius_} and \code{other_radius_}, updates
the height offset \code{dz_}, and then performs a normal EKF update from the
new face.

This is why the tracker state includes the enemy center and radius instead of
only tracking the visible armor point. The center lets the code predict all
four plates.

\section{Aim Planning}

The aim planner predicts four possible future armor faces:

\[
\psi_i(t) = \psi + \dot{\psi}t + i\frac{\pi}{2},
\qquad i \in \{0,1,2,3\}.
\]

The enemy center is also predicted:

\[
x_c(t) = x_c + v_{xc}t,\qquad
y_c(t) = y_c + v_{yc}t,\qquad
z_c(t) = z_a + v_{za}t.
\]

For face $i$, the face center is:

\[
x_i = x_c(t) - r_i\cos\psi_i(t),
\qquad
y_i = y_c(t) - r_i\sin\psi_i(t),
\qquad
z_i = z_c(t) + dz_i.
\]

The planner evaluates all four faces and chooses the best one. The first pass
uses a rough prediction time:

\[
t_{pred} = t_{bias} + \frac{1.5}{v_{bullet}},
\]

where $t_{bias}$ is the \param{time_bias} parameter, or measured latency when
that feature is enabled. The second pass reuses the best face's computed
bullet flight time.

\section{Barrel Origin And Parallax}

The aim vector starts at the barrel, not at the camera. This matters at close
range. A $10$ cm camera-to-barrel offset at $1$ m is a large angular error.

The barrel offset is measured in the gimbal body frame:

\[
\vect{o}_{barrel} =
\begin{bmatrix}
o_x\\o_y\\o_z
\end{bmatrix}.
\]

The code rotates horizontal offset by current yaw:

\[
b_x = o_x\cos\theta - o_y\sin\theta,
\qquad
b_y = o_x\sin\theta + o_y\cos\theta,
\qquad
b_z = h_{gimbal} + o_z.
\]

Then the aim vector is face position minus barrel position.

\section{Ballistics}

Without gravity, pitch would be:

\[
\phi = \operatorname{atan2}(\Delta z, d),
\]

where $d$ is horizontal distance and $\Delta z$ is vertical target difference.
Gravity makes the bullet drop during flight, so the barrel must aim higher.

The solver iterates:

\[
v_x = v\cos\phi,
\qquad
t = \frac{d}{v_x},
\]

\[
\phi = \operatorname{atan2}\left(\Delta z + \frac{1}{2}gt^2, d\right).
\]

After five iterations, it returns pitch and flight time. It rejects impossible
geometry such as nearly vertical shots or a pitch magnitude above about
$1.2$ radians.

\section{Fire Margin}

For each predicted face, the planner computes an angular window:

\[
w = w_0 \min\left(\frac{d_{ref}}{\max(R,0.5)},2.0\right).
\]

Here $w_0$ is \param{angular_window}, $d_{ref}$ is
\param{window_ref_dist}, and $R$ is target range.

The error is:

\[
e = |\operatorname{shortest\_angular\_distance}(bearing, face\_yaw)|.
\]

The margin is:

\[
margin = w - e.
\]

If the margin is positive, the selected face is aligned enough according to
the planner. The fire gate still performs more checks before the actual fire
flag is published.

\section{Anti-Gyro Timing}

When anti-gyro mode is enabled and the enemy yaw rate is high, face alignment
timing matters. The residual is:

\[
residual = \frac{
\operatorname{shortest\_angular\_distance}(face\_yaw,bearing)
}{\dot{\psi}}.
\]

Its unit is seconds. A residual near zero means the bullet impact time and the
face alignment time match. The fire gate can block shots when
$|residual|$ is larger than \param{fire.anti_gyro_max_residual}.

\section{Command Smoothing}

The node smooths yaw and pitch commands using an exponential moving average.
Pitch uses:

\[
p_{smooth} = \alpha p_{target} + (1-\alpha)p_{old}.
\]

Yaw uses the shortest angular distance so that a target crossing from
$+\pi$ to $-\pi$ does not command a full revolution:

\[
y_{smooth} =
normalize(y_{old} + \alpha \cdot shortest(y_{old}, y_{target})).
\]

\param{cmd_smooth_alpha}=1 means no smoothing. Smaller values smooth more but
increase lag. The fire gate blocks shooting when smoothing lag is too large.

\section{Fire Gate Logic}

The fire gate is intentionally conservative. It allows fire only when all
required checks pass:

\begin{enumerate}
  \item Tracker is in \code{TRACKING}.
  \item Aim target is valid.
  \item Ballistic solution is valid.
  \item Distance is within \param{min_fire_dist} and \param{max_fire_dist}.
  \item Optional stale-measurement and anti-gyro checks pass.
  \item Planner margin is non-negative.
  \item Camera-frame aim angle is inside \param{angular_window}.
  \item Smoothed command is close enough to the unsmoothed target.
\end{enumerate}

If a margin-like check fails briefly after a shot was allowed, hysteresis can
keep fire enabled across small angular drift. Hysteresis is dropped as soon as
a hard check fails.

\section{Debugging With The Math}

When the robot misses or jitters, debug in the same order as the math pipeline:

\begin{enumerate}
  \item Detection: Are bounding boxes stable in pixels?
  \item PnP: Is reprojection error low and depth stable?
  \item Transform: Is odom position stable when the robot is still?
  \item EKF: Are innovation and Mahalanobis distance reasonable?
  \item Aim: Is the selected face stable and the predicted target plausible?
  \item Command: Is smoothing lag large?
  \item Fire gate: Which blocker reason dominates?
\end{enumerate}

Do not start by changing EKF noise if PnP depth is already noisy. Do not tune
ballistics if the measured bullet speed is unknown. Fix the earliest wrong
stage.

\section{Practice Problems}

\begin{enumerate}
  \item A detector reports a box centered at $(500,300)$ with width $100$ px,
  height $80$ px, and \param{light_ratio}=0.85. Compute the four synthetic PnP
  image points.
  \item A point has camera coordinates $(0.2, -0.1, 3.0)$ m. With
  $f_x=f_y=700$, $c_x=480$, $c_y=300$, compute the projected pixel.
  \item If \param{bullet_speed} is accidentally set too high, will long-range
  shots tend to hit high or low? Explain using the gravity equation.
  \item A target has $\operatorname{shortest}(bearing, face\_yaw)=0.06$ rad and
  the scaled window is $0.05$ rad. What is the margin? Can the planner allow
  fire?
  \item Why is Mahalanobis distance better than plain distance for target
  association?
\end{enumerate}

\chapter{Coding Course}

\section{Goal Of This Course}

This chapter teaches the minimum coding knowledge needed to read, run, and
debug this package. It does not try to teach all of C++ or all of ROS 2. It
teaches the ideas that appear in this repository:

\begin{itemize}
  \item source files and header files,
  \item variables, functions, structs, and classes,
  \item loops and conditions,
  \item ROS 2 nodes, topics, messages, publishers, and subscribers,
  \item CMake and package metadata,
  \item practical debugging habits.
\end{itemize}

\section{What A Program Is}

A program is a set of instructions for the computer. Source code is the human
readable text that describes those instructions. A compiler translates C++
source code into machine code that the computer can run.

This package is a ROS 2 package named \repo. ROS 2 is a robotics framework. It
lets separate programs communicate by sending messages on named channels called
topics.

\section{Files, Folders, And Builds}

The important hand-written code is under \file{src/}. Generated files under
\file{build/} and \file{install/} are build output. New members should usually
read and edit \file{src/}, not generated output.

\begin{longtable}{p{0.28\textwidth}p{0.62\textwidth}}
\toprule
Folder or file & Beginner explanation \\
\midrule
\file{src/src/*.cpp} & C++ implementation files. These contain function
bodies: what the code actually does. \\
\file{src/include/auto_aim/*.hpp} & C++ header files. These declare the
types and functions other files can use. \\
\file{src/msg/*.msg} & ROS message definitions. ROS generates C++ code from
these during the build. \\
\file{src/config/*.yaml} & Parameter values loaded at runtime. This is where
robot-specific tuning belongs. \\
\file{src/launch/*.py} & Launch scripts that start the node with parameters. \\
\file{src/CMakeLists.txt} & Build instructions for CMake and colcon. \\
\file{src/package.xml} & ROS package metadata and dependencies. \\
\bottomrule
\end{longtable}

Typical build and run commands from the workspace root:

\begin{lstlisting}
colcon build --packages-select auto_aim
source install/setup.bash
ros2 launch auto_aim auto_aim_realsense_16.launch.py
\end{lstlisting}

If you edit C++ or message files, rebuild. If you only edit YAML parameters,
you usually only need to relaunch.

\section{C++ Source And Header Files}

C++ projects often split code into headers and source files.

\begin{itemize}
  \item A header file, such as \file{tracker.hpp}, says what a class or
  function looks like.
  \item A source file, such as \file{tracker.cpp}, contains the implementation.
\end{itemize}

Example from the tracker header:

\begin{lstlisting}[language=C++]
class Tracker
{
public:
  explicit Tracker(const TrackerConfig & cfg);
  void update(const std::vector<ArmorDetection> & detections, double dt);
  AimResult computeAim(double current_yaw, double current_pitch,
                       double override_pred_lead_s = -1.0) const;
};
\end{lstlisting}

This says a \code{Tracker} has a constructor, an \code{update} function, and a
\code{computeAim} function. The detailed math lives in \file{tracker.cpp}.

\section{Includes}

At the top of a C++ file, \texttt{\#include} brings in declarations from another
file or library:

\begin{lstlisting}[language=C++]
#include <rclcpp/rclcpp.hpp>
#include "auto_aim/tracker.hpp"
\end{lstlisting}

Angle brackets usually mean an external library or system header. Quotes
usually mean a header from this package.

\section{Namespaces}

The code is wrapped in:

\begin{lstlisting}[language=C++]
namespace auto_aim
{
  // code here
}
\end{lstlisting}

A namespace is a named container that prevents name collisions. The class is
\code{auto_aim::Tracker}, not just \code{Tracker}, when referred to from
outside the namespace.

\section{Variables And Types}

A variable stores a value. A type tells C++ what kind of value it stores.

\begin{lstlisting}[language=C++]
double bullet_speed = 25.0;
bool tracking = false;
int confirm_frames = 3;
std::string class_id = "0";
\end{lstlisting}

Common types in this repo:

\begin{longtable}{p{0.24\textwidth}p{0.66\textwidth}}
\toprule
Type & Meaning \\
\midrule
\code{double} & Floating point number, used for math. \\
\code{float} & Smaller floating point number, often used in ROS messages. \\
\code{int} & Whole number. \\
\code{bool} & True or false. \\
\code{std::string} & Text. \\
\code{std::vector<T>} & Resizable list of values of type \code{T}. \\
\code{std::array<T,N>} & Fixed-size list with exactly \code{N} values. \\
\code{Eigen::VectorXd} & Math vector from the Eigen library. \\
\code{cv::Mat} & Matrix/image container from OpenCV. \\
\bottomrule
\end{longtable}

\section{Structs}

A struct groups related values under one name. For example:

\begin{lstlisting}[language=C++]
struct ArmorDetection
{
  double x, y, z;
  double yaw;
  std::string class_id;
  double confidence;
};
\end{lstlisting}

Instead of passing six separate values everywhere, the code passes one
\code{ArmorDetection}.

Important structs in this package:

\begin{itemize}
  \item \code{TrackerConfig}: all tracker, ballistic, and gate parameters.
  \item \code{ArmorDetection}: one transformed detection in odom frame.
  \item \code{AimResult}: output from the aim planner.
  \item \code{DebugFrame}: internal per-frame debug data before publishing.
  \item \code{PnPResult}: output from the PnP solver.
\end{itemize}

\section{Functions}

A function is reusable code with a name. It may take inputs and return an
output.

\begin{lstlisting}[language=C++]
double Tracker::targetRange() const
{
  return std::sqrt(x_(0)*x_(0) + x_(2)*x_(2) + x_(4)*x_(4));
}
\end{lstlisting}

This function returns a \code{double}. The \code{const} at the end means it
does not modify the \code{Tracker} object.

\section{Classes}

A class groups data and functions together. In this repo, each major pipeline
piece is a class or stateless helper:

\begin{longtable}{p{0.24\textwidth}p{0.66\textwidth}}
\toprule
Class or helper & Responsibility \\
\midrule
\code{AutoAimNode} & ROS wiring: parameters, subscriptions, callbacks,
publishers. \\
\code{PnPSolver} & 2D box to 3D armor pose. \\
\code{FrameTransformer} & Camera frame to odom frame. \\
\code{Tracker} & EKF, target association, aiming, and ballistics. \\
\code{FireGate} & Fire/hold decision and reason. \\
\code{DetectionAdapter} & Convert detector messages into candidates. \\
\code{DebugPublisher} & Publish structured debug messages. \\
\code{ConfigValidator} & Validate parameters at startup. \\
\bottomrule
\end{longtable}

\section{Public And Private}

Classes often have \code{public} and \code{private} sections.

\begin{itemize}
  \item \code{public}: other code is allowed to call these functions or read
  these members.
  \item \code{private}: internal implementation details. Other code should not
  depend on them.
\end{itemize}

For example, other code can call \code{tracker_->update(...)}, but it cannot
directly change \code{x_}, the EKF state. That keeps the tracker responsible
for its own consistency.

\section{Pointers And Unique Ownership}

Some objects are stored through pointers:

\begin{lstlisting}[language=C++]
std::unique_ptr<Tracker> tracker_;
tracker_ = std::make_unique<Tracker>(cfg);
\end{lstlisting}

A pointer stores the location of an object. \code{std::unique_ptr} means one
owner is responsible for that object. When the owner is destroyed, the object
is automatically cleaned up.

To call a function through a pointer, use \code{->}:

\begin{lstlisting}[language=C++]
tracker_->update(armors, dt);
\end{lstlisting}

\section{References And \texorpdfstring{\code{const}}{const}}

You will often see parameters like:

\begin{lstlisting}[language=C++]
void update(const std::vector<ArmorDetection> & detections, double dt);
\end{lstlisting}

\code{&} means pass by reference. The function can use the original object
without copying the whole vector. \code{const} means the function promises not
to modify that vector.

This is common for performance and safety.

\section{Conditions}

\code{if} chooses between code paths:

\begin{lstlisting}[language=C++]
if (!pnp_.ready()) return;
if (!imu_valid_) {
  RCLCPP_WARN_THROTTLE(get_logger(), *get_clock(), 2000,
    "Waiting for /micro_imu");
  return;
}
\end{lstlisting}

The \code{return} exits the current function early. This pattern is called a
guard clause. It keeps the rest of the function from running when required
inputs are missing.

\section{Loops}

A loop repeats work. This package often loops over detections or four armor
faces:

\begin{lstlisting}[language=C++]
for (const auto & det : detections) {
  // use det
}
\end{lstlisting}

\code{const auto & det} means: for each detection, call it \code{det}, do not
copy it, and do not modify it.

\section{ROS 2 Nodes}

A ROS 2 node is a running component that communicates with other nodes. The
\code{AutoAimNode} class inherits from \code{rclcpp::Node}:

\begin{lstlisting}[language=C++]
class AutoAimNode : public rclcpp::Node
\end{lstlisting}

This gives it ROS features such as parameters, subscriptions, publishers, and
logging.

\section{Topics And Messages}

A topic is a named channel. A message type defines what data is sent on that
channel.

\begin{longtable}{p{0.32\textwidth}p{0.32\textwidth}p{0.26\textwidth}}
\toprule
Topic & Message type & Direction \\
\midrule
\topic{/detector/armors} & \code{vision_msgs/Detection2DArray} & input \\
\topic{/camera_info} & \code{sensor_msgs/CameraInfo} & input \\
\topic{/micro_imu} & \code{std_msgs/Float32MultiArray} & input \\
\topic{/tracker/cmd_gimbal} & \code{geometry_msgs/Twist} & output \\
\topic{/cmd_vel_AI} & \code{geometry_msgs/Twist} & output legacy \\
\topic{/tracker/aim_pixels} & \code{geometry_msgs/Twist} & output debug \\
\topic{/tracker/marker} & \code{visualization_msgs/MarkerArray} & output debug \\
\topic{/auto_aim/debug} & \code{auto_aim/AutoAimDebug} & output debug \\
\bottomrule
\end{longtable}

\section{Subscriptions}

A subscription tells ROS: when a message arrives on this topic, call this
function.

\begin{lstlisting}[language=C++]
det_sub_ = create_subscription<vision_msgs::msg::Detection2DArray>(
  "/detector/armors", rclcpp::SensorDataQoS(),
  std::bind(&AutoAimNode::detectionCallback, this, std::placeholders::_1));
\end{lstlisting}

This means \code{detectionCallback} runs once for each incoming detection
frame.

\section{Publishers}

A publisher sends messages:

\begin{lstlisting}[language=C++]
cmd_pub_ = create_publisher<geometry_msgs::msg::Twist>(
  "/tracker/cmd_gimbal", rclcpp::SensorDataQoS());

geometry_msgs::msg::Twist cmd;
cmd.angular.z = yaw_target_micro;
cmd.angular.y = pitch_target_micro;
cmd.angular.x = fire_decision ? 1.0 : 0.0;
cmd_pub_->publish(cmd);
\end{lstlisting}

The firmware contract depends on these exact fields. Do not move yaw, pitch,
fire, or distance to different fields without a matching firmware change.

\section{Parameters}

Parameters are runtime settings. The node declares them at startup:

\begin{lstlisting}[language=C++]
cfg.bullet_speed = declare_parameter("bullet_speed", 25.0);
\end{lstlisting}

If a YAML file provides \param{bullet_speed}, ROS uses the YAML value. If not,
the default \code{25.0} is used.

Robot-specific values belong in \file{src/config/params_realsense_16.yaml} or
\file{src/config/params_zed_64.yaml}, not hard-coded into C++.

\section{Launch Files}

Launch files start nodes and provide parameters. For example,
\file{auto_aim_realsense_16.launch.py} loads
\file{config/params_realsense_16.yaml}. Use launch files for normal robot runs
because they keep the startup command consistent.

\section{Logging}

ROS logging prints useful runtime information:

\begin{lstlisting}[language=C++]
RCLCPP_INFO(get_logger(), "Auto-aim node started");
RCLCPP_WARN_THROTTLE(get_logger(), *get_clock(), 1000, "PnP FAILED");
\end{lstlisting}

Throttled logs print at most once per time window. This prevents the terminal
from being flooded at 30 frames per second.

\section{External Libraries Used Here}

\begin{longtable}{p{0.22\textwidth}p{0.68\textwidth}}
\toprule
Library & Why this package uses it \\
\midrule
ROS 2 \code{rclcpp} & Nodes, parameters, publishers, subscribers, logging. \\
OpenCV & PnP, camera matrices, Rodrigues rotations, point projection. \\
Eigen & Linear algebra for EKF state, matrices, covariance, and Kalman math. \\
\code{tf2} & Quaternions, vector rotations, marker orientations. \\
\code{angles} & Correct wrap-around math for radians near $+\pi$ and $-\pi$. \\
\bottomrule
\end{longtable}

\section{How To Read A New File}

Use this routine when opening a file for the first time:

\begin{enumerate}
  \item Read the matching header first if one exists.
  \item Identify the main data types: structs and classes.
  \item Find the public functions. These are the file's interface.
  \item Find where those functions are called with \code{rg "functionName"}.
  \item Read implementation from the top-level function inward.
  \item Track units: meters, radians, pixels, seconds.
  \item Track frame: camera, gimbal body, or odom.
\end{enumerate}

\section{How To Make A Small Change Safely}

\begin{enumerate}
  \item State the behavior you want to change in one sentence.
  \item Find the earliest pipeline stage responsible for that behavior.
  \item Change only that module when possible.
  \item Rebuild with \code{colcon build --packages-select auto_aim}.
  \item Relaunch and record debug data.
  \item Compare \topic{/auto_aim/debug} before and after.
\end{enumerate}

Avoid changing several unrelated parameters at once. If the behavior changes,
you will not know which change caused it.

\section{Useful Terminal Commands}

\begin{longtable}{p{0.40\textwidth}p{0.50\textwidth}}
\toprule
Command & Purpose \\
\midrule
\code{rg "text" src} & Search the source tree quickly. \\
\code{rg --files src} & List source files. \\
\code{colcon build --packages-select auto_aim} & Build only this package. \\
\code{source install/setup.bash} & Load the built package into the current shell. \\
\code{ros2 topic list} & Show active ROS topics. \\
\code{ros2 topic echo /auto_aim/debug} & Print debug messages. \\
\code{ros2 topic hz /detector/armors} & Measure detector publish rate. \\
\code{ros2 param list /auto_aim} & Show parameters for the node. \\
\code{ros2 bag record ...} & Record data for offline debugging. \\
\bottomrule
\end{longtable}

\section{Common Beginner Mistakes}

\begin{itemize}
  \item Editing generated files under \file{build/} or \file{install/}. Edit
  \file{src/} instead.
  \item Forgetting to rebuild after editing C++ or message files.
  \item Forgetting to run \code{source install/setup.bash} after a build.
  \item Mixing degrees and radians.
  \item Mixing camera-frame and odom-frame coordinates.
  \item Hiding calibration mistakes with \param{pitch_offset_deg} or
  \param{yaw_offset_deg}.
  \item Changing the \topic{/tracker/cmd_gimbal} field mapping without a
  firmware update.
\end{itemize}

\section{First Exercises}

\begin{enumerate}
  \item Find where \param{bullet_speed} is declared in the C++ node.
  \item Find where \topic{/tracker/cmd_gimbal} is created and where messages
  are published on it.
  \item Find all files that mention \code{FireGate}.
  \item Change only a YAML value, relaunch, and confirm the parameter changed
  with \code{ros2 param get}.
  \item Record \topic{/auto_aim/debug} while pointing at a static target and
  identify which fields belong to detection, PnP, EKF, aim, command, and fire
  gate.
\end{enumerate}

\chapter{Code Walkthrough}

\section{End-To-End Pipeline}

The runtime path starts when the detector publishes a
\code{vision_msgs/Detection2DArray} on \topic{/detector/armors}. The node does
the following work in one detection callback:

\begin{enumerate}
  \item Confirm camera intrinsics and IMU yaw/pitch have been received.
  \item Convert raw detections into filtered bounding-box candidates.
  \item Run PnP for each candidate.
  \item Transform each successful PnP pose from camera frame to odom frame.
  \item Reject impossible positions by height and distance.
  \item Update the tracker with the remaining armor detections.
  \item Compute an aim result from the tracker state.
  \item Evaluate the fire gate.
  \item Publish the gimbal command, legacy command, HUD pixels, RViz markers,
  and structured debug message.
\end{enumerate}

The key design decision is that the node publishes absolute yaw and pitch in
the microcontroller reference frame. It does not publish relative increments.

\section{\file{src/src/auto_aim_node.cpp}}

\subsection{Responsibility}

\file{auto_aim_node.cpp} is the ROS-facing coordinator. It owns:

\begin{itemize}
  \item all runtime parameters,
  \item subscribers for detections, camera info, and IMU yaw/pitch,
  \item publishers for commands, HUD pixels, markers, and debug data,
  \item the \code{Tracker}, \code{PnPSolver}, \code{FrameTransformer},
  \code{FireGate}, and \code{DebugPublisher} objects,
  \item command smoothing and final sign conversion.
\end{itemize}

It should stay mostly orchestration code. Heavy math belongs in
\file{tracker.cpp}, \file{pnp_solver.cpp}, or \file{frame_transformer.cpp}.

\subsection{Constructor}

The constructor declares parameters and creates all ROS interfaces. Parameter
declaration has two purposes:

\begin{enumerate}
  \item It gives every setting a default value.
  \item It allows a YAML launch file to override that value.
\end{enumerate}

The first large block fills \code{TrackerConfig}. That config is passed into
the tracker constructor:

\begin{lstlisting}[language=C++]
tracker_ = std::make_unique<Tracker>(cfg);
cfg_ = cfg;
\end{lstlisting}

Then the constructor builds a matching \code{FireGate::Config}. Fire-related
parameters are split because the tracker computes planner feasibility while
the fire gate owns the final fire/hold decision.

The target class list is converted into a \code{std::set<std::string>} for
fast lookup. Class \code{"1"} is erased because it represents grey/dead
detections and should not be tracked.

Before starting subscriptions, the constructor validates configuration with
\code{ConfigValidator::validate}. Errors abort startup. Warnings are printed
but allow the node to run. This is deliberate: a wrong sign, impossible bullet
speed, or empty target list is more dangerous than a startup failure.

\subsection{Camera Info Subscription}

The \topic{/camera_info} callback initializes the PnP solver once:

\begin{lstlisting}[language=C++]
std::array<double, 9> K;
std::copy(msg->k.begin(), msg->k.end(), K.begin());
pnp_.init(K, msg->d);
\end{lstlisting}

The matrix $K$ contains focal lengths and principal point. The distortion
vector \code{d} corrects lens distortion. The same callback stores
\code{cam_fx_}, \code{cam_fy_}, \code{cam_cx_}, and \code{cam_cy_} for HUD
pixel projection later.

\subsection{Micro IMU Callback}

\code{microImuCallback} receives \topic{/micro_imu}. The message must contain:

\[
data[0] = yaw,\qquad data[1] = pitch.
\]

Both are in radians and in the microcontroller convention. The node multiplies
them by the configured signs:

\begin{lstlisting}[language=C++]
imu_yaw_   = yaw_sign_ * micro_yaw_rad;
imu_pitch_ = pitch_sign_ * micro_pitch_rad;
\end{lstlisting}

Then it updates the frame transformer. This creates the cached rotation and
translation used by detection callbacks. If no IMU message has arrived, the
detection callback refuses to run because camera points cannot be transformed
into the correct frame.

\subsection{Detection Callback}

\code{detectionCallback} is the main hot path. It starts with two guards:

\begin{itemize}
  \item PnP must be initialized from camera info.
  \item IMU yaw/pitch must be valid.
\end{itemize}

It creates a \code{DebugFrame} at the start and fills it stage by stage. This
is why the debug topic can show where a frame failed.

Frame time \code{dt} is computed from message timestamps and clamped to
$0.01$ to $0.10$ seconds. The clamp prevents one bad timestamp from exploding
the EKF prediction.

The detection adapter filters raw detections:

\begin{lstlisting}[language=C++]
auto candidates = DetectionAdapter::convert(*msg, target_classes_);
\end{lstlisting}

For each candidate, the node runs PnP:

\begin{lstlisting}[language=C++]
auto pnp = pnp_.solve(cand.cx, cand.cy, cand.w, cand.h,
                      cfg_.light_ratio, 10.0, enable_pnp_refine_);
\end{lstlisting}

Failed PnP results are still written into the debug frame when possible. That
helps diagnose whether a detection was dropped because of PnP quality rather
than tracker logic.

Successful PnP results go through the frame transformer:

\begin{lstlisting}[language=C++]
auto td = frame_transformer_.apply(pnp.rvec, pnp.tvec, max_oblique_rad);
\end{lstlisting}

Then the node rejects positions with excessive height or range:

\begin{itemize}
  \item \param{max_armor_z} limits vertical position.
  \item \param{max_armor_distance} limits horizontal range.
\end{itemize}

The remaining detections are converted into \code{ArmorDetection} values and
passed to the tracker:

\begin{lstlisting}[language=C++]
tracker_->update(armors, dt);
\end{lstlisting}

After the tracker update, the node copies EKF state, innovation, Mahalanobis
distance, and target id into the debug frame.

\subsection{Aim Computation And Latency}

The node calls:

\begin{lstlisting}[language=C++]
auto aim = tracker_->computeAim(imu_yaw_, imu_pitch_, override_pred_lead);
\end{lstlisting}

Normally \code{override_pred_lead} is \code{-1.0}, which tells the tracker to
use the legacy \param{time_bias} calculation. If \param{use_measured_latency}
is enabled and initialized, the node passes an EMA of measured pipeline latency
plus \param{gimbal_response_delay_s}.

The callback also projects the selected impact point back into camera-relative
yaw and pitch. The fire gate uses this camera-frame total angle as the
\code{OFF_AXIS} check.

\subsection{Command Publishing}

\code{publishCommand} is where the final command is formed. The order matters:

\begin{enumerate}
  \item Start from \code{aim.abs_yaw} and \code{aim.abs_pitch}.
  \item Subtract bore-sight offsets.
  \item Store pre-smoothing values for debug.
  \item Apply EMA smoothing if tracking.
  \item Compute smoothing lag.
  \item Call \code{FireGate::evaluate}.
  \item Convert yaw and pitch back to the microcontroller convention using
  \param{gimbal.yaw_sign} and \param{gimbal.pitch_sign}.
  \item Publish \code{geometry_msgs::Twist}.
\end{enumerate}

The primary command contract is:

\begin{longtable}{p{0.24\textwidth}p{0.62\textwidth}}
\toprule
Field & Meaning \\
\midrule
\code{angular.z} & absolute yaw target in radians \\
\code{angular.y} & absolute pitch target in radians \\
\code{angular.x} & fire flag, $1.0$ means fire and $0.0$ means hold \\
\code{linear.x} & range to selected armor face in meters \\
\bottomrule
\end{longtable}

The same yaw, pitch, and fire values are also published on the legacy topic
\topic{/cmd_vel_AI}. The legacy topic does not include distance.

\subsection{Markers}

\code{publishMarkers} publishes RViz markers:

\begin{itemize}
  \item enemy center sphere,
  \item four armor cubes,
  \item velocity arrow,
  \item approximate aim point,
  \item selected impact point.
\end{itemize}

When the tracker is not tracking, the function publishes delete actions for
old markers so RViz does not show stale geometry.

\section{\file{src/include/auto_aim/debug_frame.hpp}}

\code{DebugFrame} is an internal plain data struct. The detection callback
fills it during processing. \code{DebugPublisher} copies it into the ROS
message at the end.

This two-step design keeps the hot callback readable. Each stage can write its
own fields without constructing a ROS message directly.

\section{\file{src/src/detection_adapter.cpp}}

\subsection{Responsibility}

\code{DetectionAdapter::convert} turns a
\code{vision_msgs::msg::Detection2DArray} into a vector of
\code{DetectionCandidate} structs.

It performs three checks:

\begin{enumerate}
  \item Ignore detections with no classification result.
  \item Ignore detections whose class id is not in \param{target_classes}.
  \item Ignore boxes smaller than \code{min_pixel_size}.
\end{enumerate}

The adapter does not know about PnP, tracking, or firing. That separation is
useful: detector-message parsing can be tested and changed without touching
the EKF.

\section{\file{src/src/pnp_solver.cpp}}

\subsection{Initialization}

\code{PnPSolver::init} stores camera intrinsics and distortion. It also builds
two 3D armor models:

\begin{itemize}
  \item small armor,
  \item large armor.
\end{itemize}

The object points lie on a flat plane with $x=0$ and dimensions converted from
millimeters to meters.

\subsection{Solving}

\code{PnPSolver::solve} first rejects invalid input:

\begin{itemize}
  \item non-finite center or size,
  \item boxes below minimum size,
  \item invalid \param{light_ratio}.
\end{itemize}

It then builds four synthetic image points from the bounding box. The helper
\code{try_model} runs OpenCV \code{solvePnP} with \code{SOLVEPNP_IPPE}, which
is suited for planar targets. If \param{enable_pnp_refine} is true, it runs a
Levenberg-Marquardt refinement step.

Both small and large armor models are tested. The lower reprojection error
wins. The output \code{PnPResult} contains:

\begin{itemize}
  \item \code{ok}: true only if solve succeeded and error is below threshold,
  \item \code{solver_ok}: true if PnP itself produced a finite result,
  \item \code{is_large}: which model won,
  \item \code{rvec} and \code{tvec}: OpenCV pose output,
  \item \code{reproj_err} and \code{reproj_err_norm}: quality metrics,
  \item \code{image_points}: synthetic corners used for debugging.
\end{itemize}

\section{\file{src/src/frame_transformer.cpp}}

\subsection{updateFromImu}

\code{updateFromImu} converts current yaw and pitch into a camera-to-odom
rotation.

\begin{lstlisting}[language=C++]
q_gimbal.setRPY(0.0, imu_pitch, imu_yaw);
q_conv.setRPY(-M_PI / 2.0, 0.0, -M_PI / 2.0);
rotation_ = q_gimbal * q_conv;
translation_ = tf2::Vector3(0, 0, gimbal_height_);
\end{lstlisting}

\code{q_conv} is the fixed rotation from gimbal body axes to camera optical
axes. \code{q_gimbal} changes with the IMU. The translation lifts the camera
to gimbal height.

\subsection{apply}

\code{apply} receives OpenCV \code{rvec} and \code{tvec}. It:

\begin{enumerate}
  \item converts \code{tvec} into a camera-frame point,
  \item converts \code{rvec} into a quaternion,
  \item rotates and translates position into odom,
  \item rotates orientation into odom,
  \item extracts yaw,
  \item corrects or clamps impossible armor yaw,
  \item returns \code{TransformedDetection}.
\end{enumerate}

If orientation extraction fails, \code{valid} remains false and the detection
is ignored.

\section{\file{src/src/tracker.cpp}}

\subsection{Responsibility}

\code{Tracker} is the largest math module. It owns:

\begin{itemize}
  \item EKF state \code{x_},
  \item covariance \code{P_},
  \item target state machine,
  \item target association,
  \item armor face jump handling,
  \item radius and alternate-radius tracking,
  \item ballistic aim computation,
  \item optional adaptive process noise and anti-gyro timing.
\end{itemize}

\subsection{Constructor}

The constructor initializes a 9-element zero state and a diagonal initial
covariance \code{P0_}. Larger covariance means less certainty. For example,
yaw rate starts with relatively large uncertainty because it is hard to know
from one detection.

\subsection{ekfPredict}

\code{ekfPredict} runs before measurement association unless the tracker is
\code{LOST}. It applies damped constant velocity to position, height, and yaw.

When state is \code{TEMP_LOST}, it uses \param{alpha_coast} to decay velocity
faster because no measurement is currently confirming the motion.

If \param{enable_adaptive_q} is true, process noise is changed dynamically:

\begin{itemize}
  \item stationary targets get reduced process noise to suppress jitter,
  \item high yaw-rate targets get increased yaw process noise for response.
\end{itemize}

After prediction, covariance is symmetrized and bounded. Bounding prevents
coasting uncertainty from growing without limit.

\subsection{ekfUpdate}

\code{ekfUpdate} applies a measurement
\[
\vect{z} = [x_a, y_a, z_a, yaw]^T.
\]

It builds the Jacobian \code{H}, the range- and obliquity-dependent
measurement noise \code{R}, then computes innovation, Kalman gain, state
update, and Joseph-form covariance update.

Two debug values are stored:

\begin{itemize}
  \item \code{last_innovation_norm_},
  \item \code{last_mahalanobis_}.
\end{itemize}

These are published so operators can tell whether the tracker agrees with
incoming detections.

\subsection{ekfMahalanobis}

This function computes the association distance for a candidate detection. A
candidate must have distance below \param{maha_threshold} to be associated
with the active target.

If the math solve fails, it returns a very large number so the detection is
rejected.

\subsection{initFromDetection}

When the tracker is \code{LOST}, the closest valid detection initializes the
state. The detected armor is converted into an estimated enemy center by
adding radius in the face-yaw direction:

\[
x_c = x_a + r\cos\psi,\qquad y_c = y_a + r\sin\psi.
\]

Velocities start at zero. The target id becomes the detection's class id.

\subsection{shouldSwitch}

Target switching is conservative. The tracker does not switch while
\code{DETECTING} or \code{TRACKING}. In \code{TEMP_LOST}, it can switch after a
cooldown if a new target is clearly closer according to
\param{switch_range_ratio}.

This avoids bouncing between two similar targets.

\subsection{handleArmorJump}

When the target rotates and another face becomes visible, measured yaw can
jump. \code{handleArmorJump} snaps yaw, handles possible spin reversal, updates
alternate face radius/height information, and either resets or inflates
covariance depending on whether the inferred armor position is plausible.

Then it runs a normal EKF update using the jump measurement.

\subsection{update}

\code{update} is the main per-frame tracker entry point. It has these stages:

\begin{enumerate}
  \item Decrement target-switch cooldown.
  \item Possibly switch targets if state allows it.
  \item Initialize from closest detection if state is \code{LOST}.
  \item Predict EKF state.
  \item Associate detections using class id and Mahalanobis gate.
  \item Run normal update, face-jump handling, or miss handling.
  \item Clamp radius and yaw rate.
  \item Advance the state machine.
\end{enumerate}

The state machine logic at the end is important:

\begin{itemize}
  \item \code{DETECTING} needs \param{confirm_frames} consecutive matches.
  \item \code{TRACKING} becomes \code{TEMP_LOST} after one miss.
  \item \code{TEMP_LOST} becomes \code{LOST} after \param{lost_timeout}.
\end{itemize}

\subsection{solveBallistic}

\code{solveBallistic} returns pitch and flight time for a target at horizontal
distance \code{gd} and vertical difference \code{dz}. It iterates five times
to compensate gravity. It rejects too-close, too-slow-horizontal, or extreme
pitch solutions.

\subsection{computeAim}

\code{computeAim} turns the EKF state into a gimbal command. It returns an
empty \code{AimResult} unless the tracker is \code{TRACKING} or
\code{TEMP_LOST}.

The function computes the barrel position from current yaw and configured
barrel offset. It then evaluates four future faces over two passes. Each face
candidate stores:

\begin{itemize}
  \item predicted 3D face position,
  \item range and bearing from barrel,
  \item ballistic pitch and flight time,
  \item fire margin,
  \item visibility score,
  \item anti-gyro residual when applicable.
\end{itemize}

The best face is the one with highest margin, with visibility as a tie-break.
In anti-gyro mode, the best viable residual can override the margin pick.

The final \code{AimResult} contains:

\begin{itemize}
  \item absolute yaw and pitch,
  \item relative yaw and pitch for debug/fire only,
  \item distance,
  \item selected target point,
  \item tracking and fire booleans,
  \item selected face index,
  \item flight time and prediction time,
  \item anti-gyro residual.
\end{itemize}

\section{\file{src/src/fire_gate.cpp}}

\subsection{Responsibility}

\code{FireGate} is the single source of truth for fire/hold. It receives
\code{FireGate::Inputs} from the command publisher and returns a
\code{Decision} with:

\begin{itemize}
  \item \code{fire}: true or false,
  \item \code{blocker}: enum reason,
  \item \code{reason}: human-readable debug string.
\end{itemize}

\subsection{Decision Order}

The order is deliberate:

\begin{enumerate}
  \item Not tracking blocks immediately.
  \item Invalid target blocks immediately.
  \item Invalid ballistic solution blocks immediately.
  \item Out-of-range blocks immediately.
  \item Optional stale measurement blocks.
  \item Optional anti-gyro timing blocks.
  \item Margin, off-axis, and smoothing checks decide normal fire.
  \item Hysteresis may keep fire true across small angular drift.
  \item The most informative remaining blocker is returned.
\end{enumerate}

Hard failures reset hysteresis. Hysteresis only helps small geometry-like
drifts after fire was already allowed.

\section{\file{src/src/debug_publisher.cpp}}

\code{DebugPublisher} creates the \topic{/auto_aim/debug} publisher. Its
\code{publish} function copies every \code{DebugFrame} field into the generated
ROS message type.

It also maintains a fire blocker histogram. Every
\code{fire_log_period_s} seconds it logs counts and percentages for recent
blocker reasons. This is useful when fire is not allowed and individual
messages are too noisy to inspect by eye.

\section{\file{src/src/config_validator.cpp}}

\code{ConfigValidator} checks parameter ranges before the node fully starts.
It separates warnings from errors:

\begin{itemize}
  \item Errors mean the node should not run. Examples: empty target classes,
  impossible sign values, or max fire distance less than min fire distance.
  \item Warnings mean the value is unusual but not necessarily impossible.
  Examples: nonzero bore-sight offsets or suspicious physical dimensions.
\end{itemize}

Use the validator as the first place to add checks when adding new parameters.
Bad parameter values should fail loudly at startup, not silently mis-aim on
the field.

\section{Headers In \file{src/include/auto_aim}}

The headers expose the package's internal module boundaries:

\begin{longtable}{p{0.34\textwidth}p{0.56\textwidth}}
\toprule
Header & What to read there \\
\midrule
\file{tracker.hpp} & Main data structs, tracker config, public tracker API,
state enum, and private helper declarations. \\
\file{pnp_solver.hpp} & PnP result fields, armor dimensions, solve signature. \\
\file{frame_transformer.hpp} & Frame definitions and transformed detection
output. \\
\file{fire_gate.hpp} & Fire blocker enum, input fields, config fields, and
decision output. \\
\file{detection_adapter.hpp} & Detection candidate representation and convert
function. \\
\file{debug_frame.hpp} & Internal debug fields populated across the callback. \\
\file{debug_publisher.hpp} & Debug publisher interface. \\
\file{config_validator.hpp} & Validation result and validator interface. \\
\bottomrule
\end{longtable}

\section{Messages}

\subsection{\file{src/msg/AutoAimDebug.msg}}

This is the most important message for debugging. It groups fields by pipeline
stage:

\begin{itemize}
  \item detection,
  \item PnP,
  \item frame transform,
  \item EKF,
  \item aim planner,
  \item command publisher,
  \item fire gate,
  \item latency.
\end{itemize}

When a bug appears, inspect these groups in order. For example, if
\code{pnp_tvec_z} is unstable on a static target, do not tune the EKF first.

\subsection{\file{src/msg/GimbalCmd.msg}}

This message is not used at runtime. The actual command topic is
\topic{/tracker/cmd_gimbal} with \code{geometry_msgs/Twist}. Keep
\file{GimbalCmd.msg} only as a placeholder unless the team intentionally
migrates the firmware contract.

\section{Configuration Files}

\file{params_realsense_16.yaml} and \file{params_zed_64.yaml} provide
robot-specific runtime values. The two files are similar, but physical values
such as \param{gimbal_height}, \param{barrel_offset_x}, and
\param{barrel_offset_z} should be measured per robot.

Important categories:

\begin{itemize}
  \item detection/tracking: target classes, light ratio, max armor range,
  state machine settings;
  \item EKF: process noise, measurement noise, damping;
  \item physical: bullet speed, gravity, gimbal height, barrel offset;
  \item fire gate: angular window, range limits;
  \item timing: prediction bias, reference frequency;
  \item command shaping: smoothing alpha and sign convention.
\end{itemize}

If the code has a default value and the YAML has a different value, the YAML
wins at launch.

\section{Launch Files}

\file{auto_aim_realsense_16.launch.py} and \file{auto_aim_zed_64.launch.py}
load the matching YAML file and start \code{auto_aim_node}. They assume
upstream nodes provide:

\begin{itemize}
  \item \topic{/detector/armors},
  \item \topic{/camera_info},
  \item \topic{/micro_imu}.
\end{itemize}

\file{auto_aim_launch.py} declares parameters inline. It is useful as an older
single-file example, but the robot-specific YAML launch files are clearer for
real deployments.

\section{Build Files}

\subsection{\file{src/CMakeLists.txt}}

CMake does the following:

\begin{enumerate}
  \item sets C++17,
  \item finds ROS 2, OpenCV, Eigen, tf2, and message-generation dependencies,
  \item generates message code from \file{msg/GimbalCmd.msg} and
  \file{msg/AutoAimDebug.msg},
  \item builds \code{libauto_aim_component.so} from all C++ source files,
  \item links generated message typesupport,
  \item registers \code{auto_aim::AutoAimNode} as a composable node and
  standalone executable,
  \item installs headers, launch files, config files, and docs.
\end{enumerate}

If a new C++ source file is added, it must be listed in
\code{add_library(...)} or it will not be compiled.

\subsection{\file{src/package.xml}}

\file{package.xml} is ROS package metadata. It names dependencies so ROS tools
know what must be present to build and run the package.

\section{Debugging Recipes By Symptom}

\subsection{Node Starts But Publishes No Commands}

Check:

\begin{enumerate}
  \item Is \topic{/camera_info} publishing?
  \item Is \topic{/micro_imu} publishing at least two floats?
  \item Is \topic{/detector/armors} publishing?
  \item Does \topic{/auto_aim/debug} show \code{detection_present=true}?
  \item Is tracker state stuck in \code{LOST} or \code{DETECTING}?
\end{enumerate}

\subsection{PnP Fails Often}

Check:

\begin{enumerate}
  \item Bounding boxes are not too small.
  \item \param{light_ratio} is reasonable.
  \item Camera intrinsics are correct.
  \item \code{pnp_reproj_err_norm} is not large on clean frames.
\end{enumerate}

\subsection{Target Position Jitters While Robot Is Still}

Check in order:

\begin{enumerate}
  \item bbox center and size jitter,
  \item PnP \code{tvec},
  \item \topic{/micro_imu} yaw and pitch noise,
  \item odom position,
  \item EKF innovation and covariance behavior.
\end{enumerate}

\subsection{Pitch Is Consistently Wrong}

Check:

\begin{enumerate}
  \item measured \param{bullet_speed},
  \item \param{gimbal_height},
  \item \param{barrel_offset_z},
  \item pitch sign,
  \item whether \param{pitch_offset_deg} was used to hide a calibration issue.
\end{enumerate}

\subsection{Fire Is Never Allowed}

Check \code{fire_blocker} and \code{fire_blocker_reason}. Common causes:

\begin{itemize}
  \item \code{NOT_TRACKING}: tracker has not reached \code{TRACKING};
  \item \code{OUT_OF_RANGE}: distance outside configured range;
  \item \code{MARGIN_NEGATIVE}: selected face is not aligned;
  \item \code{OFF_AXIS}: target is outside camera angular window;
  \item \code{SMOOTHING_LAG}: command smoothing has not caught up.
\end{itemize}

\section{Good First Code Tasks}

\begin{enumerate}
  \item Add a new field to \code{DebugFrame} and \code{AutoAimDebug.msg}, then
  publish it through \code{DebugPublisher}. This teaches the full debug-message
  path.
  \item Add a validator warning for a new parameter. This teaches safe startup
  checks.
  \item Add a console log for one state transition in \code{Tracker::update}.
  This teaches tracker flow without changing math.
  \item Add a small offline unit test for \code{FireGate::evaluate}. This is a
  contained logic module with no ROS subscriptions.
\end{enumerate}

\section{Rules For Future Maintainers}

\begin{itemize}
  \item Keep units visible in comments and parameter names.
  \item Keep frame conversions inside \code{FrameTransformer} or very close to
  where data enters/leaves a frame.
  \item Keep fire decision logic centralized in \code{FireGate}.
  \item Do not change the \topic{/tracker/cmd_gimbal} contract without
  firmware coordination.
  \item Prefer adding debug fields over guessing from terminal logs.
  \item When fixing a bug, identify the earliest pipeline stage where the data
  becomes wrong.
\end{itemize}


\end{document}
